\documentclass{situationist-preprint}

\title{Fission and What Matters}
\msid{PID-03-a}
\msfamily{PID-03}
\msdiscipline{Personal Identity and the Self}
\mstype{Principal manuscript}
\mskeywords{fission; personal identity; what matters in survival; Parfit; branching; double success}

\begin{document}
\maketitle

\begin{abstract}
Suppose the two hemispheres of your brain are separated and each is transplanted into a body of its own, and each resulting person has your memories, character and plans. What happens to you? We argue that fission produces two persons, each of whom has what matters in survival, and that identity is therefore not what matters. Neither successor is you in the sense of the psychological criterion of PID-01: the case branches symmetrically, so the bodily tie-breaker cannot select one, and identity cannot hold with both. Yet each successor has what matters, because if either hemisphere had been transplanted alone the successor would have had it, and whether a successor has what matters depends only on its relation to the original, not on whether another successor exists. The contrast with identity, which is extrinsic in exactly this respect (PID-02), is what makes fission informative. We defend the conclusion against the claims that fission is death, that the original survives as both, and that uniqueness is part of what matters.
\end{abstract}

\section{The problem}

The case of fission was introduced by \citet{Wiggins1967} as a difficulty for psychological accounts of personal identity and turned by \citet{Parfit1971} into an argument about what matters. Each hemisphere of a human brain can, within limits, sustain a person's psychology; this is suggested by the survival of patients after hemispherectomy and by the partial independence of the hemispheres in split-brain patients, which MND-15 interprets as partial unity. Imagine, then, that your hemispheres are separated and housed in two bodies, Lefty and Righty, each psychologically continuous with you.

There are four answers to the question what happens to you \citep[ch.~12]{Parfit1984}: you do not survive; you survive as exactly one of the two; you survive as both; or the case produces two persons, neither of whom is you in the strict sense, though each stands to you in the relations that ordinarily make survival valuable. \citet{Parfit1971,Parfit1984} defended the last answer and drew the lesson that identity is not what matters. \citet{Lewis1976} accepted the lesson's practical import while arguing that, on a perdurantist account, what matters and identity coincide after all, since the two persons share their pre-fission stages. \citet{Perry1972} and \citet{Nozick1981} gave closest-continuer and related treatments, and \citet{Johnston1997} disputed that the reductionist picture supports the revision of our concerns Parfit proposed.

We defend the fourth answer, which is the catalogue's conclusion, and argue that the step from fission to the claim that identity is not what matters is valid.

\section{Framework}\label{s:frame}

\begin{definition}\label{d:fission}
A \emph{fission case} is one in which a person's brain hemispheres are separated and each is housed in a body of its own, so that each resulting being, Lefty and Righty, is psychologically continuous with the original to the same degree, through the same kind of cause, and bodily continuous with the original to the same degree.
\end{definition}

\begin{definition}\label{d:matters}
\emph{What matters in survival} is the relation, between a person and a future being, whose holding grounds the person's prudential reasons for concern about that being's future, in the way ordinary survival grounds them.
\end{definition}

\begin{assumption}\label{a:persons}
Lefty and Righty are each a person, and they are distinct from each other.
\end{assumption}

\begin{assumption}\label{a:criterion}
Identity over time for persons is given by PID-01: psychological continuity, with bodily continuity as a tie-breaker where the process branches.
\end{assumption}

\begin{assumption}\label{a:single}
In a \emph{single transplant}, in which one hemisphere is housed in a new body and the other destroyed, the resulting person is the original person and has what matters in survival.
\end{assumption}

\begin{assumption}\label{a:intrinsic}
Whether a later being has what matters in relation to an earlier person depends only on the intrinsic character of the relation between them, not on whether some other being stands in a similar relation to the earlier person.
\end{assumption}

Assumption~\ref{a:persons} follows from the facts of the case: after the operation Lefty and Righty are in different places, have different experiences and will form different plans, and each has the capacities that make something a person. Assumption~\ref{a:single} is a consequence of Assumption~\ref{a:criterion}, since the single transplant does not branch: a person who survives with half a brain, as hemispherectomy patients do, survives. Assumption~\ref{a:intrinsic} is the substantive premise. It says that what matters, unlike identity, is not subject to the extrinsic dependence that PID-02-a establishes for identity in teletransportation, where whether the replica is the original depends on when a distant body is destroyed. PID-02-a already notes that psychological connectedness lacks this extrinsic character; Assumption~\ref{a:intrinsic} says that the same holds of what matters, whatever it turns out to be.

\section{The argument}\label{s:arg}

\begin{lemma}\label{l:neither}
In a fission case, the original person is identical with neither Lefty nor Righty.
\end{lemma}

\begin{proof}
The original cannot be identical with both, since by Assumption~\ref{a:persons} Lefty and Righty are distinct, and identity is symmetric and transitive. By Assumption~\ref{a:criterion}, since the process branches, bodily continuity decides; but by Definition~\ref{d:fission} each is bodily continuous with the original to the same degree, so the tie-breaker selects neither. Nothing else distinguishes them, so identity with one and not the other would be a difference without a ground. The original is therefore identical with neither.
\end{proof}

\begin{theorem}\label{t:two}
Fission produces two persons.
\end{theorem}

\begin{proof}
By Assumption~\ref{a:persons} there are, after the operation, two distinct persons, and by Lemma~\ref{l:neither} neither is the original.
\end{proof}

\begin{theorem}\label{t:matters}
In a fission case, each of Lefty and Righty has what matters in survival.
\end{theorem}

\begin{proof}
Consider Lefty. Compare the fission case with the single transplant in which Lefty's hemisphere is transplanted and Righty's destroyed. By Assumption~\ref{a:single}, in the single transplant Lefty has what matters. The intrinsic relation between the original and Lefty is the same in both cases: the same hemisphere, carrying the same psychology through the same cause, housed in the same kind of body. The cases differ only in whether Righty exists. By Assumption~\ref{a:intrinsic} that difference does not affect whether Lefty has what matters. So Lefty has what matters in the fission case. The argument for Righty is the same.
\end{proof}

\begin{corollary}\label{c:notid}
Identity is not what matters in survival.
\end{corollary}

\begin{proof}
By Lemma~\ref{l:neither} the original is identical with neither successor, and by Theorem~\ref{t:matters} each successor has what matters. So what matters can hold without identity, and identity is not the relation of Definition~\ref{d:matters}.
\end{proof}

Corollary~\ref{c:notid} is Parfit's point about double success: if having one successor is as good as ordinary survival, having two cannot be as bad as death \citep[ch.~12]{Parfit1984}. The argument does not show that identity has no value at all. PID-06 argues that psychological connectedness carries most of what matters and that identity matters less than commonly assumed; what Corollary~\ref{c:notid} shows is that identity is not the relation in virtue of which survival matters.

\begin{remark}\label{r:perd}
On the perdurantist account of PID-11, the two persons are sums of stages that share their pre-fission stages \citep{Lewis1976}. Theorem~\ref{t:two} then says that after fission there are two persons; Lemma~\ref{l:neither} says that the pre-fission person, counted at a time by her stage, is not uniquely either. Lewis's claim that what matters and identity coincide in this framework depends on counting the pre-fission occupant as two persons, and so on a count by sums rather than by stages. On either count Theorem~\ref{t:matters} stands, and what matters is the relation between stages, which is what Corollary~\ref{c:notid} requires.
\end{remark}

\section{Variants}\label{s:var}

\paragraph{Asymmetric fission.} Suppose Lefty receives the left hemisphere together with the original body, while Righty's hemisphere is housed in a new body. The process branches, and by Assumption~\ref{a:criterion} the tie-breaker now selects Lefty, who is the original. This is the structure of the Branch-Line case of PID-02-a. The argument of Theorem~\ref{t:matters} still applies to Righty: in the single transplant in which only Righty's hemisphere survived, Righty would have been the original and would have had what matters, and the presence of Lefty does not alter Righty's relation to the original. So Righty has what matters without being the original. Asymmetric fission therefore yields the conclusion of Corollary~\ref{c:notid} without the symmetry premise, and shows that the corollary does not depend on Lemma~\ref{l:neither}'s appeal to groundless differences.

\paragraph{Unequal hemispheres.} Suppose Lefty's hemisphere carries the language functions and most autobiographical memory, Righty's much less. Then Lefty is more strongly connected with the original than Righty. The argument shows that each has what the corresponding single transplant would give; if the single transplant of Righty's hemisphere gives only part of what matters, Righty has only that part. What matters comes in degrees, as PID-06 argues, and fission distributes it according to the connections each successor bears; it does not halve it.

\paragraph{Delayed fission.} Suppose one hemisphere is stored and transplanted a year after the other. Then there is a year in which the process has not yet branched. Once the second recipient is activated while the first exists, the process branches, and by the reasoning of PID-02-a's temporal condition the earlier recipient was not the original, though for a year no fact about it or its relation to the original revealed this. The verdict on identity changes with timing; by Assumption~\ref{a:intrinsic}, the verdict on what matters does not.

\section{Objections and replies}\label{s:obj}

\begin{objection}[Fission is death]
If you will be identical with neither successor, you will cease to exist, and ceasing to exist is as bad as death, whoever replaces you.
\end{objection}

\begin{reply}
The objection assumes that identity is what matters, which is what is in question. Theorem~\ref{t:matters} gives an argument, from the single transplant, that each successor has what matters. To resist it, the objector must say either that the single transplant does not give what matters, which no one believes of a hemispherectomy survivor, or that the existence of a second successor takes it away. The second option makes the value of your future depend on something that happens in another body and makes no difference to your relation to your successor.
\end{reply}

\begin{objection}[Survival as both]
The original survives as both successors, as a person with a divided body \citep{Lewis1976}.
\end{objection}

\begin{reply}
After fission Lefty and Righty have separate streams of consciousness and incompatible plans; a single person with two such streams would be a divided self, which PID-15 shows to involve at most partial unity, and fission involves none. The perdurantist version of the view is compatible with ours (Remark~\ref{r:perd}); it relabels the count and concedes the substance.
\end{reply}

\begin{objection}[Uniqueness matters]
Part of what matters is being the only one: the only person with my relationships, my property and my place in others' lives. Fission destroys uniqueness, so neither successor has all of what matters.
\end{objection}

\begin{reply}
Uniqueness is valuable, but it is a good of a person's external relations, like being married to a particular person or holding a particular office, and not the relation in virtue of which survival matters to the one who survives. Each successor loses some of these goods, as does a person whose identical twin is unexpectedly found. That is a loss within a life, not the loss of a life. It is part of the small residue that PID-06 allows identity and its correlates to carry.
\end{reply}

\begin{objection}[Extrinsic what matters]
Assumption~\ref{a:intrinsic} begs the question: if identity is what matters, and identity is extrinsic (PID-02-a), then what matters is extrinsic too.
\end{objection}

\begin{reply}
Assumption~\ref{a:intrinsic} is supported independently. The reasons that ground prudential concern are reasons for caring about the experiences, projects and character of a future being, and these are fixed by its relation to the earlier person: what it remembers, intends and values, and how. That a further being, causally isolated from the first, has similar features does not alter any of these. It is the extrinsic character of identity, established in PID-02-a, that is surprising; the conclusion that identity is not what matters removes the surprise.
\end{reply}

\section{What the result does not establish}\label{s:not}

The argument does not establish how much of what matters each successor has relative to ordinary survival, beyond the claim that each has what the single transplant gives; PID-06 addresses the weighting. It does not establish that the original's special obligations pass to the successors, or how. And it does not decide whether the thing that each of us fundamentally is, which PID-05 identifies as a thinking organism, survives fission; on that view the organism ceases or continues according to the fate of its life, and the persons of Theorem~\ref{t:two} are persons in the phase-sortal sense PID-05-a defines.

\section{Conclusion}

Fission produces two persons, neither of whom is the original, since the case branches symmetrically and the tie-breaker of PID-01 selects neither. Each has what matters, since each would have had it in a single transplant and the existence of the other changes nothing in its relation to the original. Identity is therefore not what matters in survival.

\begin{thebibliography}{99}\small
\bibitem[Johnston(1997)]{Johnston1997} Johnston, M. (1997). Human concerns without superlative selves. In J. Dancy (ed.), \emph{Reading Parfit}. Blackwell.
\bibitem[Lewis(1976)]{Lewis1976} Lewis, D. (1976). Survival and identity. In A.~O. Rorty (ed.), \emph{The Identities of Persons} (pp.~17--40). University of California Press.
\bibitem[Nozick(1981)]{Nozick1981} Nozick, R. (1981). \emph{Philosophical Explanations}. Harvard University Press.
\bibitem[Parfit(1971)]{Parfit1971} Parfit, D. (1971). Personal identity. \emph{The Philosophical Review}, 80(1), 3--27.
\bibitem[Parfit(1984)]{Parfit1984} Parfit, D. (1984). \emph{Reasons and Persons}. Oxford University Press.
\bibitem[Perry(1972)]{Perry1972} Perry, J. (1972). Can the self divide? \emph{The Journal of Philosophy}, 69(16), 463--488.
\bibitem[Wiggins(1967)]{Wiggins1967} Wiggins, D. (1967). \emph{Identity and Spatio-Temporal Continuity}. Blackwell.
\end{thebibliography}

\end{document}
